% =====================================================================
% Drop-in subsection for AttnFuse_paper.tex.
% Insert near the H100 results section. Requires \usepackage{booktabs}.
% =====================================================================
\subsection{Measured Hardware-Counter Profile (H100)}
\label{sec:ncu}

The analytical roofline in \S\ref{sec:roofline} states what the kernel
\emph{could} achieve; this subsection states what it actually does.
Table~\ref{tab:ncu} reports Nsight Compute counters on an H100 SXM5
(sm\_90, CUDA 12.x, driver 535) for the forward kernel of each variant
at $N{=}4096$, $B{=}4$, $H{=}12$, $D{=}64$, fp16. The harness
(\texttt{benchmarks/ncu\_profile.py}) warms up Triton's JIT cache for
three launches and profiles the fourth so the measurement is free of
one-time compile cost.

\begin{table}[h]
\centering
\small
\begin{tabular}{lrrrrr}
\toprule
\textbf{Counter} & \textbf{Dense} & \textbf{Causal} & \textbf{SW-256} & \textbf{C+ALiBi} & \textbf{RoPE+C} \\
\midrule
HMMA pipe \% peak                       & 30.5 & 15.8 & 15.2 & 15.0 & 11.5 \\
SM throughput \% peak                   & 65.7 & 35.5 & 38.1 & 37.8 & 35.5 \\
Warp occupancy \% peak                  & 24.2 & 12.0 & 12.1 & 12.0 & 12.0 \\
DRAM throughput \% peak                 &  2.0 &  2.1 &  6.6 &  2.0 &  1.5 \\
Wait stalls (math-pipe deps) \%         & 13.2 & 29.9 & 28.3 & 28.4 & 29.3 \\
Short-scoreboard stalls (SMEM) \%       &  7.7 & 19.5 & 18.5 & 19.5 & 13.4 \\
Long-scoreboard stalls (HBM) \%         &  1.4 &  1.2 &  3.8 &  1.0 & 12.8 \\
\midrule
Launch: regs/thread                     &  119 &  217 &  243 &  218 &  232 \\
Launch: SMEM/block (KB)                 &   65 &   65 &   65 &   65 &   49 \\
Launch: block size (threads)            &  256 &  128 &  128 &  128 &  128 \\
\bottomrule
\end{tabular}
\caption{Nsight Compute counters for the AttnFuse forward kernel on
an H100 SXM5 at $N{=}4096$, fp16. ``Wait stalls'' counts cycles a warp
is stalled on tensor-core math-pipe dependencies. SMEM and HBM stalls
are reported per-warp-active. All percentages are \texttt{ncu}'s
\texttt{pct\_of\_peak\_sustained\_elapsed} (or \texttt{\_active}) metric.}
\label{tab:ncu}
\end{table}

Three observations follow.

\paragraph{The H100 gap is a tensor-core feeding problem, not a
bandwidth problem.} DRAM throughput sits at 1.5--6.6\% of peak across
all five variants, and long-scoreboard stalls (the canonical HBM-bound
indicator) are below 4\% on every variant except RoPE+causal. The
kernel is firmly compute-bound on H100. Yet HMMA pipe utilisation is
only 11.5--30.5\% of peak, with 13.2--29.9\% of cycles stalled waiting
on tensor-core math-pipe dependencies. This is the diagnostic
signature of an Ampere-shaped HMMA tile pipeline running on Hopper:
each warp issues a small (\texttt{16$\times$16$\times$16}) HMMA, the
dependency latency between issues exceeds the pipeline depth, and
Hopper's larger WGMMA tiles, which would amortise the dependency
across $\sim$8$\times$ more work per instruction, are not emitted by
the current code path.

\paragraph{The Ampere sparse-table heuristic inverts on Hopper.} On
the 3090 the sparse table (causal, sliding-window, ALiBi) wins by
shrinking BLOCK\_M from 128 to 64, raising the program count and
filling more SMs. On H100 the same shrink doubles registers per
thread (119$\to$217--243), halves block size (256$\to$128), and
collapses warp occupancy to 12\%. The combined effect is roughly
2$\times$ worse HMMA utilisation and 2$\times$ worse SM throughput
than the dense path. Re-tuning the sparse Hopper table to keep
BLOCK\_M=128 (or larger with WGMMA) is the lowest-effort improvement
the profiling identifies.

\paragraph{RoPE+causal has the most distinctive profile.} It is the
only variant with material HBM-bound stalls (long-scoreboard 12.8\%,
vs.\ 1--4\% elsewhere) and the lowest HMMA utilisation (11.5\%). The
in-loop K rotation issues a gather over \texttt{rot\_offs\_d} that
accesses cos/sin tables held in SMEM, but the K tile itself is HBM-
resident on each iteration, and on Hopper the latency of that load is
no longer hidden by the (smaller) compute behind it. This is the
variant where the structural fusion is most valuable and where the
Hopper-native codegen has the largest headroom to recover.

\paragraph{Implications for Phase 1 Hopper codegen.} The profile
directly motivates the three changes already listed in the future-work
discussion (\S\ref{sec:future}): (i) emit WGMMA via Triton's
\texttt{warp\_specialize=True} path with BLOCK\_M=128, BLOCK\_N=256;
(ii) replace HBM loads of K/V with TMA descriptors
(\texttt{tl.experimental\_descriptor\_load}) so the K rotation
overlaps with copy issue; (iii) re-derive the sparse-variant tile
table from a Hopper-side sweep rather than porting the Ampere entries.
The expected effect: HMMA utilisation moving from 15--30\% toward the
60--70\% range typical of CUTLASS Hopper kernels at this shape, and
the 1.49--2.14$\times$ gap to \texttt{flex\_attention}
(Table~\ref{tab:flex_h100}) closing.
